\documentclass[11pt]{article}
\usepackage{amsmath, amssymb, amsthm}
\usepackage[margin=1in]{geometry}
\usepackage[most]{tcolorbox}
\usepackage{microtype}
\usepackage[hidelinks]{hyperref}

\definecolor{accent}{HTML}{2F5BEA}
\newtcbtheorem[number within=section]{theorem}{Theorem}%
  {colback=accent!4, colframe=accent!70, fonttitle=\bfseries, boxrule=0.5pt, arc=2pt}{thm}
\newtcbtheorem[use counter from=theorem]{definition}{Definition}%
  {colback=black!2, colframe=black!40, fonttitle=\bfseries, boxrule=0.5pt, arc=2pt}{def}
\theoremstyle{remark}
\newtheorem*{example}{Example}

\title{Linear Algebra --- Lecture 1}
\author{Your Name}
\date{\today}

\begin{document}
\maketitle
\tableofcontents

\section{Vector spaces}

\begin{definition}{Vector space}{vs}
A \emph{vector space} over a field $F$ is a set $V$ with addition and scalar multiplication satisfying the usual axioms.
\end{definition}

\begin{example}
$\mathbb{R}^n$ with componentwise operations is a vector space over $\mathbb{R}$.
\end{example}

\begin{theorem}{Rank--nullity}{rank}
For a linear map $T \colon V \to W$ with $V$ finite-dimensional,
\[
  \dim V = \operatorname{rank} T + \operatorname{nullity} T.
\]
\end{theorem}

See Theorem~\ref{thm:rank}.

\end{document}
